\documentclass[../mathNotesPreamble]{subfiles}

\providecommand{\relscalefact}{1.4}
\begin{document}
\relscale{\relscalefact}
  \section{3.7: Rational Functions}
  \begin{defn*}[Rational Function]
    A \textbf{rational function} is a function of the form
      \[r(x)=\dfrac{f(x)}{g(x)},\]
    where $f(x)$ and $g(x)$ are polynomials.
  \end{defn*}
  \vspace*{\stretch{1}}
  
  \begin{center}
    \begin{tikzpicture}[
      custNode/.style={
        black,
        align=center,
        inner sep=2.5pt,
        rounded corners,
        fill=white,
        opacity=0.8,
        text opacity=1.0},
      ]
      \begin{groupplot}[
        group style={group size=2 by 1,horizontal sep=3cm},
        grid=both, %major,minor
        grid style={line width=0.375pt, draw=gray!60},
        minor grid style={draw=gray!75},
        axis lines=center,
        axis line style={black,->},
        xmin=-6, xmax=6,
        ymin=-6, ymax=6,
        enlargelimits={abs=0.5},
        ticklabel style={font=\footnotesize,inner sep=0.5pt,fill=white,opacity=0.5, text opacity=1},
        xlabel=, ylabel=,
        every axis plot/.append style={domain=xmin:xmax, line width=0.95pt, color=lander_blue, samples=255},]
        \nextgroupplot
          \addplot[<->] expression[domain=-6:-1/6]{1/x};
          \addplot[<->] expression[domain=1/6:6]{1/x};
          \node[custNode] at (4,4) {$f(x)=\dfrac{1}{x}$};
        \nextgroupplot
          \addplot[<->] expression[domain=-6:-1/sqrt(6)]{1/x^2};
          \addplot[<->] expression[domain=1/sqrt(6):6]{1/x^2};
          \node[custNode] at (4,4) {$f(x)=\dfrac{1}{x^2}$};
      \end{groupplot}
    \end{tikzpicture}
  \end{center}
  \vspace*{\stretch{1}}
  
  %vertical asymptotes
  \begin{defn*}[Asymptotes]
    A \textbf{vertical asymptote} of a graph is a vertical line $x=a$ where the graph tends towards positive or negative infinity as the input approaches $a$ from either the left or the write. We write
      \[\textnormal{As } x \to a^-,\ f(x)\to \pm\infty \quad\textnormal{ or }\quad x\to a^+,\ f(x)\to \pm\infty.\]
    A \textbf{horizontal asymptote} of a graph is a horizontal line $y=b$ where the graph approaches the line as the inputs increase or decrease without bound. We write
      \[\textnormal{As } x\to \infty \textnormal{ or } x\to\infty,\ f(x)\to b.\]
  \end{defn*}
  \pagebreak

  %end behavior - horizontal asymptotes, slant asymptotes
  \begin{ex*}
    Describe the asymptotes, intercepts, and domain of 
      \[f(x)=\dfrac{1}{(x+1)(x-2)^2}\]
    using the graph below:
    \begin{center}
      \begin{tikzpicture}[declare function={
        f(\x)=1/((\x+1)*(\x-2)^2);
        }]
        \begin{axis}[
          grid=both, %major,minor
          grid style={line width=0.375pt, draw=gray!50},
          minor grid style={draw=gray!75},
          axis lines=center,
          axis line style={black,->},
          xmin=-6, xmax=6,
          ymin=-6, ymax=6,
          minor x tick num=1,
          enlargelimits={abs=0.5},
          ticklabel style={font=\footnotesize,inner sep=0.5pt,fill=white,opacity=0.5, text opacity=1},
          xlabel=, ylabel=,
          width=0.65\linewidth, height=1.1*\axisdefaultheight,
          every axis plot/.append style={domain=xmin:xmax, line width=0.95pt, color=lander_blue, samples=255},
          ]
          \addplot[<->] expression[domain=-6:-1.0169]{f(x)};
          \addplot[<->] expression[domain=-0.98271:1.76408]{f(x)};
          \addplot[<->] expression[domain=2.21863:6]{f(x)};
        \end{axis}
      \end{tikzpicture}
    \end{center}
  \end{ex*}
  \pagebreak

  \begin{ex*}
    Describe the asymptotes, intercepts, and domain of 
      \[f(x)=\dfrac{2x^4-2x^3+x^2+2x-2}{\parens{x^2-1}\parens{x^2+1}}\]
    using the graph below:
    \begin{center}
      \begin{tikzpicture}[declare function={
        f(\x)=(2*x^4-2*x^3+x^2+2*x-2)/((\x^2-1)*(\x^2+1));
        }]
        \begin{axis}[
          grid=both, %major,minor
          grid style={line width=0.375pt, draw=gray!50},
          minor grid style={draw=gray!75},
          axis lines=center,
          axis line style={black,->},
          xmin=-6, xmax=6,
          ymin=-6, ymax=6,
          minor x tick num=1,
          enlargelimits={abs=0.5},
          ticklabel style={font=\footnotesize,inner sep=0.5pt,fill=white,opacity=0.5, text opacity=1},
          xlabel=, ylabel=,
          width=0.65\linewidth, height=1.1*\axisdefaultheight,
          every axis plot/.append style={domain=xmin:xmax, line width=0.95pt, color=lander_blue, samples=255},
          ]
          \addplot[<->] expression[domain=-6:-1.07373]{f(x)};
          \addplot[<->] expression[domain=-0.97404:0.96724]{f(x)};
          \addplot[<->] expression[domain=1.04645:6]{f(x)};
        \end{axis}
      \end{tikzpicture}
    \end{center}
  \end{ex*}
  \pagebreak

  \begin{ex*}
    Describe the asymptotes, intercepts, and domain of 
      \[f(x)=\dfrac{x^2-3x+4}{2\parens{x-1}}\]
    using the graph below:
    \begin{center}
      \begin{tikzpicture}[declare function={
        f(\x)=(\x^2-3*\x+4)/(2*(\x-1));
        }]
        \begin{axis}[
          grid=both, %major,minor
          grid style={line width=0.375pt, draw=gray!50},
          minor grid style={draw=gray!75},
          axis lines=center,
          axis line style={black,->},
          xmin=-6, xmax=6,
          ymin=-6, ymax=6,
          minor x tick num=1,
          enlargelimits={abs=0.5},
          ticklabel style={font=\footnotesize,inner sep=0.5pt,fill=white,opacity=0.5, text opacity=1},
          xlabel=, ylabel=,
          width=0.65\linewidth, height=1.1*\axisdefaultheight,
          every axis plot/.append style={domain=xmin:xmax, line width=0.95pt, color=lander_blue, samples=255},
          ]
          \addplot[<->] expression[domain=-6:0.83095]{f(x)};
          \addplot[<->] expression[domain=1.14435:6]{f(x)};
        \end{axis}
      \end{tikzpicture}
    \end{center}
  \end{ex*}

  \pagebreak
\end{document}
